Model the ferrite bead as a series inductance L at low-mid frequencies. The two capacitors $C_1$ (input side) and $C_2$ (output side) form a pi filter with that series inductance. The pi network has a resonance frequency $f$ where the L and the effective capacitance interact:
\begin{equation}
    C_{eq} = \frac{C_1 \times C_2}{C_1 + C_2}
\end{equation}
The resonance / cutoff frequency of the $L$-$C_{eq}$ combination is:
\begin{equation}
    f = \frac{1}{2\pi \sqrt{L \times C_{eq}}}
\end{equation}
At frequencies well above $f_r$ the filter attenuates; below fr the caps and ferrite appear largely as bypasses so DC and low frequency pass through.
\nl
For ferrite beads, its resonant frequency is typically 100MHz (unless stated otherwise). At frequencies well below resonance, the bead behaves approximately like an inductor:
\begin{equation}
    Z \approx j2\pi f L
\end{equation}
Rearranging:
\begin{equation}
    L \approx \frac{Z}{2\pi f}
\end{equation}